Penrose's proposal relates the mass at spatial infinity to the area
required to enclose a black-hole region, as a consequence expected from
gravitational collapse and cosmic censorship~\cite{Penrose1973}.
In the time-symmetric setting, Huisken and Ilmanen proved the Riemannian
inequality for a connected horizon by weak inverse mean curvature
flow~\cite{HuiskenIlmanen2001}; Bray's conformal flow established the
inequality using the total area of an outermost minimal boundary with
several components~\cite{Bray2001}.
For charged Riemannian data, Khuri, Weinstein and Yamada proved the sharp
inequality with multiple horizon components under the appropriate area
branch condition, together with an upper bound for the area radius
without that condition~\cite{KhuriWeinsteinYamada2017}.

Angular momentum introduces additional boundary and gauge information.
Dain, Khuri, Weinstein and Yamada established mass--charge--angular-momentum
and lower-area bounds for simply connected, axisymmetric maximal data
with two ends~\cite{DainKhuriWeinsteinYamada2013}.
Khuri, Sokolowsky and Weinstein obtained a Penrose-type bound for maximal
axisymmetric data with a connected minimal boundary; their comparison
retains an integral of the horizon data and gives the Kerr--Newman
expression under additional horizon matching
conditions~\cite{KhuriSokolowskyWeinstein2019}.
Their harmonic-map method is a useful starting point for the comparison
developed here.

The deformation and first-variation methods draw on the manuscript
\emph{Equality and rigidity in the spacetime Penrose inequality}~\cite{NeutralExterior2026}.
The related charged companion~\cite{ConditionalCharged2026} incorporates
a numerical theorem whose earlier formulation assumed both a neutral
exterior inequality and a separate elliptic existence input; the companion
proves those inputs internally.
We establish the deformation and equality arguments needed here at
their stated scope; no earlier neutral, charged or rotating Penrose
inequality is assumed as an input.
